\section{Recherche par Similarit�}

Nous reprendrons les listes inverses cr��es dans la section~\ref{LI}, question~\ref{LI_q}.

\subsection{Produit Scalaire}

\begin{enumerate}
  \item Calculer le produit scalaire de la requ�te suivante : ``France Ralentie PIB''
\begin{correction}
\begin{itemize}
  \item doc$_1 = \frac{1}{66} \times \log{\frac{4}{3}} \times \frac{1}{3} + \frac{2}{66} \times \log{\frac{4}{3}} \times \frac{1}{3} + \frac{2}{66} \times \log{\frac{4}{2}} \times \frac{1}{3} = 4,934\times 10^{-3}$
  \item doc$_2 = \frac{1}{68} \times \log{\frac{4}{3}} \times \frac{1}{3} + \frac{1}{68} \times \log{\frac{4}{3}} \times \frac{1}{3} + \frac{0}{68} \times \log{\frac{4}{2}} \times \frac{1}{3} = 1,225\times 10^{-3}$
  \item doc$_3 = \frac{2}{97} \times \log{\frac{4}{3}} \times \frac{1}{3} + \frac{0}{97} \times \log{\frac{4}{3}} \times \frac{1}{3} + \frac{2}{97} \times \log{\frac{4}{2}} \times \frac{1}{3} = 2,928\times 10^{-3}$
  \item doc$_4 = \frac{0}{93} \times \log{\frac{4}{3}} \times \frac{1}{3} + \frac{2}{93} \times \log{\frac{4}{3}} \times \frac{1}{3} + \frac{0}{93} \times \log{\frac{4}{2}} \times \frac{1}{3} = 0,895\times 10^{-3}$
  \end{itemize}
\end{correction}
  \item Trier le r�sultat par ordre de pertinence d�croissante
\begin{correction}
$doc_1, doc_3, doc_2, doc_4$
\end{correction}
\end{enumerate}  




\subsection{Cosinus}
\label{cos}
\begin{enumerate}
  \item Donner le calcul de la norme d'un document et de la requ�te
\begin{correction}
$\sqrt{\sum tf_i^2}$
\end{correction} 
	\item Les normes des documents sont respectivement : 0.0392, 0.0336, 0.0310, 0.0421
  \item Calculer le cosinus de la requ�te suivante : ``France Ralentie PIB''
\begin{correction}
\begin{itemize}
  \item doc$_1 = \frac{4,934\times 10^{-3}}{0.0392} = 0,126$
  \item doc$_2 = \frac{1,225\times 10^{-3}}{0.0336} = 0,036$
  \item doc$_3 = \frac{2,928\times 10^{-3}}{0.0310} = 0,094$
  \item doc$_4 = \frac{0,895\times 10^{-3}}{0.0421} = 0,021$
  \end{itemize}
\end{correction}
  \item Trier le r�sultat par ordre de pertinence d�croissante
\label{cos_q}
\begin{correction}
$R_1 : doc_1, doc_3, doc_2, doc_4$
\end{correction}
\end{enumerate}

\subsection{Pr�cision / Rappel}

Nous nous attendions � l'ensemble r�sultat suivant : $R_2 : doc_1, doc_4$
\begin{enumerate}
  \item Donner la pr�cision et le rappel du r�sultat obtenu par Cosinus
\begin{correction}
\begin{itemize}
  \item $\frac{R_1 \cap R_2}{R_1} = \frac{2}{4}$
  \item $\frac{R_1 \cap R_2}{R_2} = \frac{2}{2}$
\end{itemize}
\end{correction}
  \item Donner le pr�cision et le rappel du r�sultat obtenu par Cosinus avec un score sup�rieur � 0,03
\begin{correction}
\begin{itemize}
  \item $\frac{R_1 \cap R_2}{R_1} = \frac{1}{3}$
  \item $\frac{R_1 \cap R_2}{R_2} = \frac{1}{2}$
\end{itemize}
\end{correction}
\end{enumerate}  
